% Section 8: Discussion and Limitations
% Target length: 300-500 words
% Status: OUTLINE ONLY

\section{Discussion and Limitations}
\label{sec:discussion}

% TODO: expand each into prose.
\begin{itemize}
    \item \textbf{Evidence level.} The originating observation is a subjective
          quality read across two projects, not a controlled study. Section
          \ref{sec:evaluation} is designed to address this, but until it is run
          and reported, CIS-weighted allocation should be presented as a
          well-motivated hypothesis, not a confirmed result.
    \item \textbf{Unrecoverable comparison data.} The Google Antigravity/Gemini
          side of the original Claude-vs-Gemini observation was not preserved
          (no committed learning graph or chapter content in the source repo).
          This limits the paper to a single-model (Claude) empirical pipeline;
          the cross-model claim is anecdotal and should be flagged as such, or
          reframed as future work (re-run the same 3-condition study with a
          second model).
    \item \textbf{Single domain.} Initial data is drawn from one Algebra I course
          (200 concepts). Generalization to other domains (humanities, embedded
          systems, etc., where dependency structure and "elaboration" norms may
          differ) is untested.
    \item \textbf{Normalization sensitivity.} The global-vs-local scope
          decision (\S\ref{sec:formal-definitions}) meaningfully changes which
          concepts get Tier A treatment. This is a real design choice, not a
          solved problem -- report it as such rather than presenting one option
          as obviously correct.
    \item \textbf{Discretizing a heavy-tailed recursive metric is easy to get
          wrong.} Our first tiering rule (population percentile rank) passed
          inspection on Chapter 1 alone but was structurally broken: CIS has a
          hard floor, most concepts in a real learning graph are leaves tied
          at that floor, and population rank cannot distinguish among them --
          Tier C was unreachable regardless of actual importance. This was
          only caught by validating on a second, differently-shaped chapter
          (\S\ref{sec:system}). The general point -- any recursive/PageRank-style
          importance measure with a floor and a heavy-tailed distribution needs
          value-based, not rank-based, discretization, and needs validation on
          more than one example before the discretization rule is trusted --
          likely generalizes beyond this paper's specific metric and may be
          worth stating as a standalone methodological caution for others
          applying similar graph-centrality metrics to content generation.
    \item \textbf{Cost/complexity tradeoff.} Adding a graph-analysis step
          (CIS computation) to an authoring pipeline is a small one-time cost per
          book (linear time, per Proposition 1) but is additional process
          complexity relative to a flat word-count instruction. Worth stating this
          plainly rather than assuming the more sophisticated method is always
          worth adopting.
    \item \textbf{Risk of metric gaming / overfitting structure.} A generation
          pipeline tuned specifically to maximize a CIS-weighted rubric could
          learn to satisfy the rubric's surface features (e.g. inserting more
          examples) without actual pedagogical improvement. The human spot-check
          in \S\ref{sec:evaluation} is a partial mitigation, not a full answer.
\end{itemize}
